% titel: Blatt 9, A6: Massenwirkungsgesetz allgemein
% Teil II: der durchdachte Weg.

\subsection*{Aufgabenstellung}
Massenwirkungsgesetz zu $aA + bB \rightleftharpoons cC + dD$ aufstellen.

\subsection*{a) Massenwirkungsgesetz}
\[
  K_c = \frac{[C]^c\,[D]^d}{[A]^a\,[B]^b}.
\]
Koeffizienten $a,b,c,d$ werden zu den Exponenten.

\subsection*{b) Ändert sich $K_c$ bei Verdopplung von [A]?}
Nein. $K_c$ ist eine Konstante bei fester Temperatur. Wird [A] verdoppelt,
wächst der Quotient über $K_c$; nach Le Chatelier reagiert das System,
bis wieder $Q=K_c$ gilt --- aber der \emph{Wert von $K_c$ bleibt gleich}.
Nur die Gleichgewichtslage (Mengen) verschiebt sich, die Konstante nicht.

\subsection*{c) Wie kann man $K_c$ verändern?}
Durch \textbf{Änderung der Temperatur}. Nur die Temperatur beeinflusst die
Gleichgewichtskonstante (Arrhenius-/Van-'t-Hoff-Abhängigkeit). Die
Veränderung von Konzentrationen, Druck oder Volumen verschiebt das
Gleichgewicht, lässt $K_c$ aber unverändert. Ein Katalysator ändert $K_c$
ebenso wenig.

\begin{ergebnis}
$K_c = \frac{[C]^c[D]^d}{[A]^a[B]^b}$; Verdopplung von [A] ändert $K_c$
nicht; $K_c$ ändert sich nur über die Temperatur.
\end{ergebnis}

\begin{kontrolle}
Konsistent mit der Definition: $K$ hängt nur von $T$ ab. Fragen zur
``Änderung von $K_c$`` müssen daher immer auf die Temperatur abzielen.
\end{kontrolle}

\begin{merksatz}
$K_c$ ist temperatur-abhängig und sonst konstant. Konzentrationsänderungen
verschieben nur die Lage, nie $K_c$ selbst.
\end{merksatz}
